\Needspace{6\baselineskip}
\begin{center}\hypertarget{native-259-article186}{}\pdfbookmark[2]{186.}{article-186}186.\end{center}
Ceterum principium, quod quadrata differentiarum inter quantitates observatas
et computatas summam quam minimam producere debeant, etiam independenter
a calculo probabilitatis sequenti modo considerari poterit.

Quoties multitudo incognitarum multitudini quantitatum observatarum
independentium aequalis est, illas ita determinare licet, ut his exacte satisfiat.
Quoties autem multitudo illa hac minor est, consensus absolute exactus obtineri
nequit, quatenus observationes praecisione absoluta non gaudent. In hoc
itaque casu operam dare oportet, ut consensus quam optimus stabiliatur, sive
ut differentiae quantum fieri potest extenuentur. Haec vero notio natura sua
aliquid vagi involvit. Etiamsi enim systema valorum pro incognitis, quod
\emph{omnes} differentias resp. minores reddit quam aliud, procul dubio huic praeferendum
sit, nihilominus optio inter duo systemata, quorum alterum in aliis
observationibus consensum meliorem offert, alterum in aliis, arbitrio nostro
quodammodo relinquitur, manifestoque innumera principia diversa proponi
possunt, per quae conditio prior impletur. Designando differentias inter observationes
et calculum per $\Delta$, $\Delta'$, $\Delta''$ etc., conditioni priori non modo satisfiet,
si $\Delta\Delta+\Delta'\Delta'+\Delta''\Delta''+\text{etc.}$ fit minimum (quod est principium nostrum), sed etiam
si $\Delta^4+\Delta'^4+\Delta''^4+\text{etc.}$, vel $\Delta^6+\Delta'^6+\Delta''^6+\text{etc.}$, vel generaliter summa potestatum
exponentis cuiuscunque paris in minimum abit. Sed ex omnibus his
principiis nostrum simplicissimum est, dum in reliquis ad calculos complicatissimos
deferremur. Ceterum principium nostrum, quo iam inde ab anno 1795
usi sumus, nuper etiam a clar. \textsc{Legendre} in opere \emph{Nouvelles méthodes pour la
\SourceLeaf{260}détermination des orbites des comètes, Paris} 1806, prolatum est, ubi plures aliae
proprietates huius principii expositae sunt, quas hic brevitatis caussa supprimimus.

Si potestatem exponentis paris infinite magni adoptaremus, ad systema id
reduceremur, in quo differentiae maximae fiunt quam minimae.

Ill. \textsc{Laplace} ad solutionem aequationum linearium, quarum multitudo maior
est quam multitudo quantitatum incognitarum, principio alio utitur, quod olim
iam a clar. \textsc{Boscovich} propositum erat, scilicet ut differentiae ipsae sed omnes
positive sumtae summam minimam conficiant. Facile ostendi potest, systema
valorum incognitarum, quod ex hoc solo principio erutum sit, necessario\footnote{Casibus specialibus exceptis, ubi solutio quodammodo indeterminata manet.}
tot aequationibus e propositarum numero exacte satisfacere debere, quot sint
incognitae, ita ut reliquae aequationes eatenus tantum in considerationem veniant,
quatenus \emph{ad optionem decidendam conferunt}: si itaque e.\,g. aequatio
$V=M$ est ex earum numero, quibus non satisfit, systema valorum secundum
illud principium inventorum nihil mutaretur, etiamsi loco ipsius $M$ valor quicunque
alius $N$ observatus esset, si modo designando per $n$ valorem computatum
differentiae $M-n$, et $N-n$ eodem signo affectae sint. Ceterum ill.
\textsc{Laplace} principium istud per adiectionem conditionis novae quodammodo
temperat: postulat scilicet, ut summa differentiarum ipsa, signis non mutatis,
fiat $=0$. Hinc efficitur, ut multitudo aequationum exacte repraesentatarum
unitate minor fiat quam multitudo quantitatum incognitarum, verumtamen
quod ante observavimus etiamnum locum habebit, siquidem duae saltem incognitae
affuerint.
